%======================================================================
%  What a cluster is, and which part of one this thesis uses
%======================================================================

\section{Supercomputers, clusters and nodes}
\label{sec:hpc-primer}

\subsection{What makes a machine a supercomputer}

A modern supercomputer is not a single very large computer.  It is a
few hundred to a few thousand ordinary computers, wired together and
managed as one facility.  Each of those computers, called a \textbf{node}, has its own
processors, its own memory and its own operating system, and looks from
the inside much like a well-specified server.  What turns
a room full of nodes into a supercomputer is the three things bolted
around them: a fast \textbf{interconnect} so that nodes can exchange
data, a \textbf{parallel file system} they all see, and a
\textbf{scheduler} that hands nodes out to users a job at a time.

\figref{fig:workstation-vs-cluster} contrasts this with the machine on
a desk.  The important difference is not size but \emph{addressability}.
On a workstation, every core can read every byte of memory directly.  On
a cluster that remains true \emph{within} a node and stops being true
between nodes: to use another node's memory, a program has to send a
message over the network and wait.  Crossing that boundary is roughly
three orders of magnitude more expensive than a local memory access, and
almost everything difficult about parallel programming follows from it.

\begin{figure}[htbp]
\centering
\begin{tikzpicture}[
  font=\small,
  box/.style={draw=black!60, fill=neutralcol, rounded corners=2pt,
              minimum width=2.5cm, minimum height=0.62cm, align=center,
              font=\footnotesize},
  nodebox/.style={draw=black!55, fill=white, rounded corners=2pt,
                  minimum width=1.5cm, minimum height=1.15cm,
                  align=center, font=\scriptsize},
  svc/.style={draw=ddrcol!80!black, fill=ddrcol!10, rounded corners=2pt,
              minimum width=2.3cm, minimum height=0.62cm, align=center,
              font=\footnotesize},
  lbl/.style={font=\footnotesize, align=center},
  lnk/.style={-{Stealth}, thick, black!60}
]
  % workstation side
  \node[font=\bfseries] at (-4.4,3.05) {A workstation};
  \node[box] (wc) at (-4.4,2.25) {CPU cores};
  \node[box, fill=hbmcol!14, draw=hbmcol!80!black] (wm) at (-4.4,1.35)
    {memory};
  \node[box] (wd) at (-4.4,0.45) {disk};
  \draw[lnk] (wc) -- (wm);
  \draw[lnk] (wm) -- (wd);
  \node[lbl, text width=4.4cm, black!75] at (-4.4,-0.95)
    {One operating system.\\\emph{Every core can address\\every byte
     directly.}};

  \draw[black!25, dashed] (-1.55,-1.95) -- (-1.55,3.3);

  % cluster side
  \node[font=\bfseries] at (4.3,3.05) {A cluster};
  \node[svc] (login) at (1.35,2.25) {login node};
  \node[svc] (sched) at (1.35,1.35) {scheduler};
  \draw[lnk] (login) -- (sched);

  \foreach \i/\x in {1/3.6, 2/5.1, 3/6.6} {
    \node[nodebox] (n\i) at (\x,2.05)
      {\textbf{node}\\{\tiny CPUs}\\{\tiny memory}};
  }
  \node[font=\scriptsize, black!60] at (7.75,2.05) {$\cdots$};

  \draw[lnk] (sched) -- (2.85,2.05);
  \draw[black!60, thick] (3.0,0.95) -- (7.9,0.95);
  \foreach \x in {3.6,5.1,6.6} \draw[black!60, thick] (\x,1.47) -- (\x,0.95);
  \node[lbl, black!70] at (5.45,0.62) {interconnect};

  \node[svc, minimum width=4.9cm] (fs) at (5.45,-0.15)
    {parallel file system};
  \draw[black!60, thick] (5.45,0.95) -- (5.45,0.16);

  \node[lbl, text width=6.6cm, black!75] at (4.9,-1.32)
    {Many operating systems.\\\emph{A core addresses its own node's
     memory directly;\\anything else is a message over the network.}};
\end{tikzpicture}
\caption{A workstation and a cluster, drawn to show where direct memory
addressing stops.  The dashed line is the boundary this thesis works
inside: all of the experiments run on \emph{one} node, so the
interconnect and the scheduler are context rather than subject.
Author's synthesis.}
\label{fig:workstation-vs-cluster}
\end{figure}

\subsection{How a cluster is used}

Access follows a fixed shape, and it is worth stating because it
constrains the experimental method described in Part~III.  A user logs
in to a \textbf{login node}, which is shared, modest, and meant for
editing and compiling rather than for running anything.  Work is then
submitted to a \textbf{batch scheduler}, which on CRESCO8 is Slurm, as
a script that declares what it needs: how many nodes, how
many cores, for how long, from which partition.  The scheduler queues
the request, and when the resources free up it runs the script on the
allocated nodes and returns the output.

Two consequences matter here.  First, \emph{measurements are taken
without an interactive session}, so anything the experiment needs to
record must be written down by the job script at the time; there is no
opportunity to inspect the machine afterwards.  Second, \emph{the
allocation is not the whole machine}.  A job receives specific nodes
from a specific partition, and the properties of those nodes, rather than
the properties advertised for the cluster, are what a result depends
on.  This is why \secref{sec:first-allocation} treats capturing a
machine-readable system manifest as the first task of the first
allocation rather than as housekeeping.

\subsection{Why HPC hardware differs from a desktop}

Nodes in a cluster are built for sustained throughput rather than for
peak responsiveness, and the differences show up in exactly the places
this thesis cares about.  They carry two or more processor sockets, so
memory access time depends on which socket a thread runs on.  They
carry many more memory channels than a desktop, because feeding fifty
or a hundred cores takes considerably more bandwidth than feeding
eight.  They are cooled aggressively, on CRESCO8 with water, so that high core
counts can be sustained rather than briefly reached.  And
increasingly they carry \emph{more than one kind of memory}, which is
where this thesis begins.

That last point is the reason a cluster appears in a study of language
model inference at all.  The processor examined here, the Intel Xeon CPU
Max~9480, is a mainstream server CPU with an unusual property: it has
high-bandwidth memory in the package alongside ordinary DDR5, and both
can be addressed by software.  Machines built from it are rare, and
essentially all of them are supercomputers.  Studying the question
therefore requires cluster access even though the question itself is
about a single socket.

\subsection{What this thesis actually uses}

The funnel narrows quickly, and \tabref{tab:hpc-scope} states where it
stops.  Of everything above, the experiments use one node from one
partition, and the primary comparison uses one socket of that node.  The
interconnect is never on the critical path.  The parallel file system
matters only for storing model files.  The scheduler matters only
because it decides which node is obtained and for how long.

\begin{table}[htbp]
\centering
\small
\tabsetup
\caption{Which parts of the cluster this thesis depends on.  The
narrowness is deliberate: every level that is excluded is a level whose
variability cannot contaminate the memory-tier comparison.}
\label{tab:hpc-scope}
\begin{tabularx}{\textwidth}{@{}P{3.05cm}P{3.3cm}Y@{}}
\toprule
\headrow
\textbf{Level} & \textbf{Role in this thesis} & \textbf{Why} \\
\midrule
Cluster
& \textbf{Context only}
& Provides access to a processor that is otherwise unobtainable, and
  imposes the queue limits that shape the experiment matrix. \\
\addlinespace[1pt]
Interconnect
& \textbf{Excluded}
& All experiments are single-node, so network performance never enters
  a reported measurement. \\
\addlinespace[1pt]
File system
& \textbf{Support only}
& Holds model files; excluded from every timing boundary. \\
\addlinespace[1pt]
Scheduler
& \textbf{Method}
& Determines which node is allocated, which is why the observed node
  configuration, not the documented one, is what gets reported. \\
\addlinespace[1pt]
Node
& \textbf{Secondary subject}
& Two sockets joined by UPI; used for the topology questions, where
  crossing between sockets is the effect being measured. \\
\addlinespace[1pt]
\textbf{Socket}
& \textbf{Primary subject}
& \emph{The unit on which the central comparison is defined.}  One set
  of cores, one local HBM tier, one local DDR5 tier. \\
\bottomrule
\end{tabularx}
\end{table}

With that scope fixed, the next two sections work inwards:
\secref{sec:xeon-max-platform} describes the socket and the memory
system that make the question askable, and \secref{sec:cresco8}
describes the particular cluster on which it will be asked.
